% Figure: the scrape-off-layer flux tube striking the divertor at a shallow
% angle. The parallel heat flux carried along the field is enormous, but the
% field meets the target at a grazing angle so the flux is spread over a
% larger footprint, cutting the surface load by the sine of the angle. The
% sketch shows the last closed flux surface, the narrow scrape-off layer of
% width lambda_q, and the tilted target that spreads the load.
\begin{figure}[htbp]
\centering
\begin{tikzpicture}[>=Latex, scale=1.0]
  % last closed flux surface (plasma edge)
  \draw[engblue, thick] (0,0.6) .. controls (2.5,1.4) and (5.0,1.2) .. (7.0,0.4);
  \node[engblue, font=\scriptsize, anchor=south] at (2.0,1.25)
    {last closed flux surface};
  % scrape-off layer, a thin band just outside
  \draw[engorange, thick] (0,0.35) .. controls (2.5,1.15) and (5.0,0.95) .. (7.0,0.15);
  \draw[<->, enggray] (7.15,0.15) -- (7.15,0.4);
  \node[enggray, font=\scriptsize, anchor=west] at (7.2,0.28)
    {$\lambda_q$};
  \node[engorange, font=\scriptsize, anchor=south] at (5.4,0.95)
    {scrape-off layer};
  % field line spiralling down into the divertor along the SOL
  \draw[engred, thick, ->] (6.6,0.25) .. controls (7.2,-0.4) and (7.0,-1.1) .. (6.2,-1.7);
  \node[engred, font=\scriptsize, anchor=west] at (7.0,-0.8)
    {$q_\parallel$};
  % divertor target plate, tilted at a shallow angle
  \draw[engblack, very thick] (4.3,-2.3) -- (7.7,-1.3);
  \node[engblack, font=\scriptsize, anchor=north] at (6.0,-2.0)
    {divertor target};
  % grazing angle arc
  \draw[enggray] (6.2,-1.7) -- (6.9,-2.15);
  \draw[enggray, ->] (6.7,-1.55) arc (250:290:0.6);
  \node[enggray, font=\scriptsize] at (6.55,-1.35) {$\alpha$};
  % footprint bracket on the plate
  \draw[decorate, decoration={brace, amplitude=4pt}, enggray]
    (5.9,-2.35) -- (7.1,-1.98);
  \node[enggray, font=\scriptsize, anchor=north] at (6.5,-2.45)
    {spread footprint};
\end{tikzpicture}
\caption[Divertor heat load and the grazing target]{The scrape-off layer
carries the plasma exhaust along the field to the divertor target. The heat
flux parallel to the field, $q_\parallel$, is enormous because the power that
crosses the last closed flux surface is funnelled into a scrape-off layer only
a few millimetres wide, $\lambda_q$. The target is set at a shallow grazing
angle $\alpha$ to the field so that the parallel flux is projected onto a
larger surface footprint, and the load the material sees is
$q_\parallel\sin\alpha$, smaller than the parallel flux by the sine of the
grazing angle. Tilting the target and widening the wetted footprint are the
two geometric knobs a divertor designer has; when they are not enough, the
exhaust power is radiated away volumetrically before it reaches the plate by
seeding an impurity that lights up in the scrape-off layer, the detached
divertor regime.}
\label{fig:vol04ch13:divertor-flux}
\end{figure}
